%!TEX root = ../template.tex
%%%%%%%%%%%%%%%%%%%%%%%%%%%%%%%%%%%%%%%%%%%%%%%%%%%%%%%%%%%%%%%%%%%%
%% chapter3-attacks-draft.tex
%%
%% DRAFT continuation of Chapter 3 (Methodology), intended to follow the
%% (A)/(B) itemised list at the end of Section~\ref{subsec:pgd_maddpg}.
%% Standalone: not yet \input by 0-Config/4_files.tex and chapter3.tex is
%% untouched. Copy in, or add %!TEX root = ../template.tex
%%%%%%%%%%%%%%%%%%%%%%%%%%%%%%%%%%%%%%%%%%%%%%%%%%%%%%%%%%%%%%%%%%%%
%% chapter3-attacks-draft.tex
%%
%% DRAFT continuation of Chapter 3 (Methodology), intended to follow the
%% (A)/(B) itemised list at the end of Section~\ref{subsec:pgd_maddpg}.
%% Standalone: not yet \input by 0-Config/4_files.tex and chapter3.tex is
%% untouched. Copy in, or add %!TEX root = ../template.tex
%%%%%%%%%%%%%%%%%%%%%%%%%%%%%%%%%%%%%%%%%%%%%%%%%%%%%%%%%%%%%%%%%%%%
%% chapter3-attacks-draft.tex
%%
%% DRAFT continuation of Chapter 3 (Methodology), intended to follow the
%% (A)/(B) itemised list at the end of Section~\ref{subsec:pgd_maddpg}.
%% Standalone: not yet \input by 0-Config/4_files.tex and chapter3.tex is
%% untouched. Copy in, or add %!TEX root = ../template.tex
%%%%%%%%%%%%%%%%%%%%%%%%%%%%%%%%%%%%%%%%%%%%%%%%%%%%%%%%%%%%%%%%%%%%
%% chapter3-attacks-draft.tex
%%
%% DRAFT continuation of Chapter 3 (Methodology), intended to follow the
%% (A)/(B) itemised list at the end of Section~\ref{subsec:pgd_maddpg}.
%% Standalone: not yet \input by 0-Config/4_files.tex and chapter3.tex is
%% untouched. Copy in, or add \input{2-MainMatter/chapter3-attacks-draft}.
%%%%%%%%%%%%%%%%%%%%%%%%%%%%%%%%%%%%%%%%%%%%%%%%%%%%%%%%%%%%%%%%%%%%

% ─────────────────────────────────────────────────────────────────────────────
\subsection{Threat model and admissible perturbation set}
\label{subsec:threat_model_pgd}

Both extensions share a single threat model, which is deliberately identical to the
one used for the FGSM baseline of Section~\ref{subsec:fgsm_maddpg_def} so that any
observed difference in damage is attributable to the attack itself rather than to a
change in the attacker's capabilities.

Let $\mathcal{C} \subseteq \{1,\dots,M\}$ denote the set of compromised agents and
$N = |\mathcal{C}|$ its cardinality. The victim policies $\{\pi_{\theta_i}\}$ and
critics $\{Q_{\phi^i}\}$ are \emph{frozen}: their parameters are fixed throughout,
and the adversary interacts with them only through forward evaluations and
gradients taken with respect to the \emph{input} observations. No victim parameter
is ever updated, so the attack measures the robustness of a fixed trained policy
rather than the outcome of an adversarial training game.

The adversary perturbs observations only; the underlying network state evolves from
the true, unperturbed dynamics. For agent $i$ at time $t$ the admissible set is
\begin{equation}
  \mathcal{S}\big(s_{i,t}^{\text{local}}, \varepsilon\big) =
  \Big\{ \hat{s} \;\Big|\;
    \big\| \hat{s} - s_{i,t}^{\text{local}} \big\|_{\infty} \le \varepsilon,
    \;\; \hat{s} \in \mathcal{D}
  \Big\},
  \label{eq:admissible_set}
\end{equation}

\noindent where $\mathcal{D}$ is the domain constraint that re-projects the
normalised bandwidth features onto $[0,1]$. The corresponding projection operator
$\Pi_{\mathcal{S}}$ therefore composes an $\ell_{\infty}$ clip about the clean
observation with the domain clip of Equation~\eqref{eq:fgsm_maddpg_obs}. Because
the measured damage depends on the size of the admissible set, the \emph{same}
operator $\Pi_{\mathcal{S}}$ is applied to every attack and to the random control
of Section~\ref{subsec:eval_protocol}; using a stricter domain clip for one attack
than another would confound the perturbation direction with the size of the set it
is drawn from.

% ─────────────────────────────────────────────────────────────────────────────
\subsection{(A) Coordinated multi-agent perturbation}
\label{subsec:coordinated_attack}

\subsubsection*{Separability of the per-agent attack}

The FGSM baseline, and any attack that iterates over agents independently, solves
$N$ decoupled problems: for each $i \in \mathcal{C}$,
\begin{equation}
  \eta_{i,t}^{\star} = \arg\max_{\eta \,:\, s_{i,t}^{\text{local}} + \eta \in \mathcal{S}}
  \; \mathcal{L}_{i}\big(s_{i,t}^{\text{local}} + \eta\big),
  \label{eq:separable_attack}
\end{equation}

\noindent where $\mathcal{L}_i$ depends only on agent $i$'s own observation and
policy. The joint perturbation obtained this way is the concatenation of $N$
independent solutions, and the implied joint objective is \emph{separable},
\begin{equation}
  \mathcal{J}_{\text{sep}}(\hat{\mathbf{s}}_t)
  = \sum_{i \in \mathcal{C}} \mathcal{L}_i(\hat{s}_{i,t}),
  \label{eq:separable_objective}
\end{equation}

\noindent so no term couples two agents. This is a structural limitation rather
than a matter of attack strength: a separable objective cannot express a preference
for outcomes that depend on \emph{which} agents are steered onto the \emph{same}
resource, because the value of steering agent $i$ onto a link is independent of
whether agent $j$ was steered onto it as well. In a topology whose robustness comes
from $K$-path redundancy, the damaging configurations are precisely those in which
several redirected flows converge on one surviving link, and these lie outside the
reach of Equation~\eqref{eq:separable_objective}.

\subsubsection*{Joint formulation}

We therefore optimise over the joint perturbation directly. Writing
$\mathbf{s}_t = \big(s_{1,t}^{\text{local}}, \dots, s_{N,t}^{\text{local}}\big)$ for
the stacked observations of the compromised agents, the coordinated attack solves
\begin{equation}
  \hat{\mathbf{s}}_t^{\star}
  = \arg\max_{\hat{\mathbf{s}} \,\in\, \mathcal{S}(\mathbf{s}_t, \varepsilon)^{N}}
  \; \mathcal{J}\big(\hat{\mathbf{s}}\big),
  \label{eq:joint_attack}
\end{equation}

\noindent with a \emph{non-separable} global objective $\mathcal{J}$. A single
gradient of $\mathcal{J}$ supplies a descent direction for every compromised agent
simultaneously, so one optimisation step trades the agents off against one another.

\subsubsection*{Differentiating through the routing decision}

The environment decodes a routing decision by taking, for each destination $d$, the
$\arg\max$ over the $K$ candidate paths. The realised route is therefore a piecewise
constant function of the observation, and its gradient vanishes almost everywhere,
which would leave Equation~\eqref{eq:joint_attack} without usable first-order
information. We adopt a straight-through estimator: the forward pass uses the hard
per-destination one-hot decision actually executed by the environment, while the
backward pass propagates the gradient of the underlying softmax policy,
\begin{equation}
  \mathrm{ST}\big(\pi_{\theta_i}(\hat{s})\big)
  = \operatorname{sg}\!\Big[
      \operatorname{onehot}\big(\pi_{\theta_i}(\hat{s})\big) - \pi_{\theta_i}(\hat{s})
    \Big] + \pi_{\theta_i}(\hat{s}),
  \label{eq:straight_through}
\end{equation}

\noindent where $\operatorname{sg}[\cdot]$ denotes the stop-gradient operator and
$\operatorname{onehot}(\cdot)$ is applied blockwise, one block of $K$ entries per
destination. This choice matters for interpreting the results: the gradient is
informative about how to \emph{change a routing decision}, which is not the same
quantity as how to \emph{cause congestion}, and the two coincide only when the
network cannot absorb the redirected flow.

\subsubsection*{Attack objectives}

Two instantiations of $\mathcal{J}$ are considered. The first is grounded in the
victim's own centralised critics, evaluated on the joint action induced by the
perturbed observations,
\begin{equation}
  \mathcal{J}_{\text{critic}}(\hat{\mathbf{s}})
  = -\frac{1}{N} \sum_{i \in \mathcal{C}}
    Q_{\phi^i}\Big( x_t, \; \big(\tilde{a}_1, \dots, \tilde{a}_N \big) \Big),
  \qquad
  \tilde{a}_j = \mathrm{ST}\big(\pi_{\theta_j}(\hat{s}_j)\big),
  \label{eq:obj_critic}
\end{equation}

\noindent where $x_t$ is the true centralised state, held fixed during the attack.
Because the centralised critic of the CTDE formulation
(Section~\ref{subsec:actor_critic}) takes the actions of \emph{all} agents as
input~\cite{lowe2017multi}, it is by construction a non-separable function of the
joint perturbation, and ascending $-Q$ steers the joint action towards the region
the victim's own value model rates as worst.

The second objective targets the shared-bottleneck mechanism explicitly. Let
$u^{(i)}_{d,k}$ denote the true bottleneck utilisation of the $k$-th candidate path
to destination $d$ for agent $i$, read from the \emph{clean} observation, and let
$[\tilde{a}_i]_{d,k}$ be the corresponding entry of the straight-through action.
Then
\begin{equation}
  \mathcal{J}_{\text{cong}}(\hat{\mathbf{s}})
  = \sum_{i \in \mathcal{C}} \sum_{d} \sum_{k}
    \big[\tilde{a}_i\big]_{d,k} \; u^{(i)}_{d,k}.
  \label{eq:obj_congestion}
\end{equation}

\noindent Summing over agents rewards configurations in which several agents are
steered onto the same congested path, which is the mechanism under test. The
utilisation coefficients are taken from the clean observation rather than the
perturbed one; were they read from $\hat{\mathbf{s}}$, the attacker would be
optimising against its own forged view of the network instead of its true state.

\subsubsection*{Joint projected gradient ascent}

Equation~\eqref{eq:joint_attack} is solved by projected gradient ascent following
the standard PGD construction~\cite{madry2019deeplearningmodelsresistant}, applied
to the stacked perturbation. Starting from a uniformly random point inside the ball,
\begin{align}
  \hat{\mathbf{s}}^{(0)}   & = \Pi_{\mathcal{S}}\big( \mathbf{s}_t + \boldsymbol{\delta}_0 \big),
    \qquad \boldsymbol{\delta}_0 \sim \mathrm{Unif}\big(\mathcal{S}(0,\varepsilon)^{N}\big), \\
  \hat{\mathbf{s}}^{(\tau+1)} & = \Pi_{\mathcal{S}}\Big(
      \hat{\mathbf{s}}^{(\tau)} + \alpha \cdot
      \operatorname{sign}\big( \nabla_{\hat{\mathbf{s}}} \mathcal{J}(\hat{\mathbf{s}}^{(\tau)}) \big)
    \Big).
  \label{eq:joint_pgd_iteration}
\end{align}

The random initialisation is what distinguishes PGD from the iterative FGSM variant,
which starts from the clean point~\cite{kurakin2017adversarialmachinelearningscale};
the procedure is repeated from $R$ independent random starts and the iterate
achieving the highest objective is retained, since the local maxima of the inner
problem are distinct but concentrated in value~\cite{madry2019deeplearningmodelsresistant}.
The step size defaults to $\alpha = 2.5\,\varepsilon / T$ for $T$ iterations, which
guarantees the boundary of the ball is reachable from any interior starting
point~\cite{madry2019deeplearningmodelsresistant}. The full procedure is given in
Algorithm~\ref{alg:coordinated_pgd}.

\begin{algorithm}[H]
\SetAlgoLined
\DontPrintSemicolon
\KwIn{clean observations $\mathbf{s}_t$; centralised state $x_t$; budget $\varepsilon$;
      steps $T$; step size $\alpha$; restarts $R$; objective $\mathcal{J}$}
\KwOut{adversarial observations $\hat{s}_{i,t}$ for all $i \in \mathcal{C}$}
$\mathcal{J}^{\text{best}} \leftarrow -\infty$\;
$U \leftarrow$ path utilisations read from the clean $\mathbf{s}_t$\;
\For{$r \leftarrow 1$ \KwTo $R$}{
  $\hat{\mathbf{s}} \leftarrow \Pi_{\mathcal{S}}\big(\mathbf{s}_t + \boldsymbol{\delta}_0\big)$,
    \quad $\boldsymbol{\delta}_0 \sim \mathrm{Unif}(\mathcal{S}(0,\varepsilon)^N)$
    \tcp*{random start}
  \For{$\tau \leftarrow 1$ \KwTo $T$}{
    $\tilde{a}_j \leftarrow \mathrm{ST}\big(\pi_{\theta_j}(\hat{s}_j)\big)$
      for all $j \in \mathcal{C}$ \tcp*{one shared graph}
    $J \leftarrow \mathcal{J}\big(\tilde{a}_{1:N}, x_t, U\big)$\;
    $g \leftarrow \nabla_{\hat{\mathbf{s}}} J$
      \tcp*{one backward $\Rightarrow$ gradient for every agent}
    $\hat{\mathbf{s}} \leftarrow \Pi_{\mathcal{S}}\big(\hat{\mathbf{s}} + \alpha \operatorname{sign}(g)\big)$\;
  }
  \If{$\mathcal{J}(\hat{\mathbf{s}}) > \mathcal{J}^{\text{best}}$}{
    $\mathcal{J}^{\text{best}} \leftarrow \mathcal{J}(\hat{\mathbf{s}})$;\quad
    $\hat{\mathbf{s}}^{\star} \leftarrow \hat{\mathbf{s}}$\;
  }
}
\Return rows of $\hat{\mathbf{s}}^{\star}$\;
\caption{Coordinated multi-agent PGD (one time step)}
\label{alg:coordinated_pgd}
\end{algorithm}

The attack is solved once per time step for the whole compromised group; each agent
subsequently receives its own block of the joint solution, so the $N$ agents act on
a single coordinated decision rather than on $N$ independent ones.

% ─────────────────────────────────────────────────────────────────────────────
\subsection{(B) Strategically-timed attacks under an \texorpdfstring{$L_0$}{L0} budget}
\label{subsec:timed_attack}

The attacks considered so far act at every time step. A more realistic adversary is
constrained in \emph{how often} it can act, which converts the attack into a
resource-allocation problem over the trajectory. We impose an $L_0$ budget: for an
episode of $T_{\text{ep}}$ steps and a budget fraction $b \in (0,1]$,
\begin{equation}
  \frac{1}{T_{\text{ep}}} \sum_{t=1}^{T_{\text{ep}}} \mathbb{1}\big[\text{attack at } t\big]
  \;\le\; b .
  \label{eq:l0_budget}
\end{equation}

Spending this budget uniformly at random would waste it on steps at which the
network has ample spare capacity. We instead spend it at moments of high leverage,
identified by a \emph{critical-state score}. We use the maximum link utilisation
across the topology,
\begin{equation}
  c_t = \max_{(u,v) \in E} \rho_t(u,v),
  \label{eq:critical_score}
\end{equation}

\noindent where $\rho_t(u,v)$ is the utilisation of link $(u,v)$ at time $t$. This
single quantity responds to both regimes of interest: congestion onset raises it
directly, and the loss of a link raises it indirectly by concentrating traffic onto
the surviving paths, so the score is elevated exactly when redundancy is thinnest.

Because the distribution of $c_t$ is not known in advance and drifts within an
episode, the trigger threshold is calibrated online. Let $W_t$ be a sliding window
of the most recent scores. The adversary acts at time $t$ if
\begin{equation}
  c_t \;\ge\; \theta_t,
  \qquad
  \theta_t = Q_{1-b}\big( \{ c_{t'} \}_{t' \in W_t} \big),
  \label{eq:timing_threshold}
\end{equation}

\noindent where $Q_{1-b}$ denotes the empirical $(1-b)$-quantile. The threshold is
computed from the window \emph{before} $c_t$ is inserted, so that a step never
influences its own decision, and the quantile construction makes the realised
attack rate track $b$ without requiring the remainder of the episode to be known in
advance. A short warm-up period, during which the adversary always acts, seeds the
window before the quantile is meaningful. The procedure is summarised in
Algorithm~\ref{alg:timing_gate}.

\begin{algorithm}[H]
\SetAlgoLined
\DontPrintSemicolon
\KwIn{budget $b$; window size $W$; warm-up length $H$; network state at time $t$}
\KwOut{decision $\textit{fire} \in \{\text{true}, \text{false}\}$}
$c_t \leftarrow \max_{(u,v) \in E} \rho_t(u,v)$ \tcp*{critical-state score}
\eIf{$|\mathcal{H}| < H$}{
  $\textit{fire} \leftarrow \text{true}$ \tcp*{warm-up: seed the window}
}{
  $\theta_t \leftarrow Q_{1-b}(\mathcal{H})$ \tcp*{computed before inserting $c_t$}
  $\textit{fire} \leftarrow (c_t \ge \theta_t)$\;
}
insert $c_t$ into $\mathcal{H}$ (ring buffer of size $W$)\;
\Return \textit{fire}\;
\caption{Critical-state timing gate}
\label{alg:timing_gate}
\end{algorithm}

The gate composes with either attack scope: it decides \emph{whether} to act at a
given step, while Section~\ref{subsec:coordinated_attack} decides \emph{how} the
perturbation is distributed across agents once the decision to act has been taken.
On skipped steps the agents receive their unmodified observations.

% ─────────────────────────────────────────────────────────────────────────────
\subsection{Evaluation protocol}
\label{subsec:eval_protocol}

\subsubsection*{Paired arms}

Each configuration is evaluated with three arms: a \emph{clean} arm with no
perturbation, a \emph{random} control that draws a uniformly random direction from
the same admissible set of Equation~\eqref{eq:admissible_set}, and the \emph{attack}
under test. All three are run on identical episodes: the traffic sequence and any
injected link failures are generated from a common seed, so episode $e$ of one arm
and episode $e$ of another differ only in the perturbation applied. The random
control draws from a dedicated random stream so that it consumes no entropy from
the traffic generator, which would otherwise desynchronise the arms and invalidate
the pairing.

\subsubsection*{Adversarial-specific damage}

Under link failures the raw degradation relative to the clean arm is dominated by
the failures themselves rather than by the attack. The primary reported quantity is
therefore the \emph{adversarial gap}, the paired per-episode difference between the
random control and the attack,
\begin{equation}
  \Delta^{\text{adv}} = \frac{1}{n} \sum_{e=1}^{n}
    \Big( \mathrm{PDR}^{\text{rand}}_{e} - \mathrm{PDR}^{\text{atk}}_{e} \Big),
  \label{eq:adversarial_gap}
\end{equation}

\noindent reported with a paired $95\%$ confidence interval,
$\Delta^{\text{adv}} \pm t_{0.975,\,n-1}\, s_d / \sqrt{n}$, where $s_d$ is the
standard deviation of the per-episode differences. Differencing within episode pairs
removes the episode-to-episode variance that otherwise dominates the comparison. The
raw drop with respect to the clean arm is reported alongside it, but is not used as
the measure of attack strength.

\subsubsection*{Decisions versus outcomes}

Two families of metric are recorded and reported separately. \emph{Decision} metrics
count how much routing behaviour the attack changed, principally the action flip
rate, the fraction of per-agent, per-destination $\arg\max$ path choices that differ
between the clean and adversarial observations at the same time step,
\begin{equation}
  \mathrm{flip}
  = \frac{1}{|\mathcal{C}|\, n_d\, T_{\text{ep}}}
    \sum_{t}\sum_{i \in \mathcal{C}}\sum_{d}
    \mathbb{1}\Big[
      \arg\max_k \pi_{\theta_i}\big(\hat{s}_{i,t}\big)_{d,k}
      \neq
      \arg\max_k \pi_{\theta_i}\big(s_{i,t}^{\text{local}}\big)_{d,k}
    \Big].
  \label{eq:flip_rate}
\end{equation}

\emph{Outcome} metrics measure delivered traffic, principally end-to-end PDR. The
distinction is essential in a topology with $K$-path redundancy, where a high flip
rate accompanied by an unchanged PDR indicates that the network absorbed the
redirected flows rather than that the attack succeeded. Reporting a flip rate as
though it were a measure of damage would conflate the two.
.
%%%%%%%%%%%%%%%%%%%%%%%%%%%%%%%%%%%%%%%%%%%%%%%%%%%%%%%%%%%%%%%%%%%%

% ─────────────────────────────────────────────────────────────────────────────
\subsection{Threat model and admissible perturbation set}
\label{subsec:threat_model_pgd}

Both extensions share a single threat model, which is deliberately identical to the
one used for the FGSM baseline of Section~\ref{subsec:fgsm_maddpg_def} so that any
observed difference in damage is attributable to the attack itself rather than to a
change in the attacker's capabilities.

Let $\mathcal{C} \subseteq \{1,\dots,M\}$ denote the set of compromised agents and
$N = |\mathcal{C}|$ its cardinality. The victim policies $\{\pi_{\theta_i}\}$ and
critics $\{Q_{\phi^i}\}$ are \emph{frozen}: their parameters are fixed throughout,
and the adversary interacts with them only through forward evaluations and
gradients taken with respect to the \emph{input} observations. No victim parameter
is ever updated, so the attack measures the robustness of a fixed trained policy
rather than the outcome of an adversarial training game.

The adversary perturbs observations only; the underlying network state evolves from
the true, unperturbed dynamics. For agent $i$ at time $t$ the admissible set is
\begin{equation}
  \mathcal{S}\big(s_{i,t}^{\text{local}}, \varepsilon\big) =
  \Big\{ \hat{s} \;\Big|\;
    \big\| \hat{s} - s_{i,t}^{\text{local}} \big\|_{\infty} \le \varepsilon,
    \;\; \hat{s} \in \mathcal{D}
  \Big\},
  \label{eq:admissible_set}
\end{equation}

\noindent where $\mathcal{D}$ is the domain constraint that re-projects the
normalised bandwidth features onto $[0,1]$. The corresponding projection operator
$\Pi_{\mathcal{S}}$ therefore composes an $\ell_{\infty}$ clip about the clean
observation with the domain clip of Equation~\eqref{eq:fgsm_maddpg_obs}. Because
the measured damage depends on the size of the admissible set, the \emph{same}
operator $\Pi_{\mathcal{S}}$ is applied to every attack and to the random control
of Section~\ref{subsec:eval_protocol}; using a stricter domain clip for one attack
than another would confound the perturbation direction with the size of the set it
is drawn from.

% ─────────────────────────────────────────────────────────────────────────────
\subsection{(A) Coordinated multi-agent perturbation}
\label{subsec:coordinated_attack}

\subsubsection*{Separability of the per-agent attack}

The FGSM baseline, and any attack that iterates over agents independently, solves
$N$ decoupled problems: for each $i \in \mathcal{C}$,
\begin{equation}
  \eta_{i,t}^{\star} = \arg\max_{\eta \,:\, s_{i,t}^{\text{local}} + \eta \in \mathcal{S}}
  \; \mathcal{L}_{i}\big(s_{i,t}^{\text{local}} + \eta\big),
  \label{eq:separable_attack}
\end{equation}

\noindent where $\mathcal{L}_i$ depends only on agent $i$'s own observation and
policy. The joint perturbation obtained this way is the concatenation of $N$
independent solutions, and the implied joint objective is \emph{separable},
\begin{equation}
  \mathcal{J}_{\text{sep}}(\hat{\mathbf{s}}_t)
  = \sum_{i \in \mathcal{C}} \mathcal{L}_i(\hat{s}_{i,t}),
  \label{eq:separable_objective}
\end{equation}

\noindent so no term couples two agents. This is a structural limitation rather
than a matter of attack strength: a separable objective cannot express a preference
for outcomes that depend on \emph{which} agents are steered onto the \emph{same}
resource, because the value of steering agent $i$ onto a link is independent of
whether agent $j$ was steered onto it as well. In a topology whose robustness comes
from $K$-path redundancy, the damaging configurations are precisely those in which
several redirected flows converge on one surviving link, and these lie outside the
reach of Equation~\eqref{eq:separable_objective}.

\subsubsection*{Joint formulation}

We therefore optimise over the joint perturbation directly. Writing
$\mathbf{s}_t = \big(s_{1,t}^{\text{local}}, \dots, s_{N,t}^{\text{local}}\big)$ for
the stacked observations of the compromised agents, the coordinated attack solves
\begin{equation}
  \hat{\mathbf{s}}_t^{\star}
  = \arg\max_{\hat{\mathbf{s}} \,\in\, \mathcal{S}(\mathbf{s}_t, \varepsilon)^{N}}
  \; \mathcal{J}\big(\hat{\mathbf{s}}\big),
  \label{eq:joint_attack}
\end{equation}

\noindent with a \emph{non-separable} global objective $\mathcal{J}$. A single
gradient of $\mathcal{J}$ supplies a descent direction for every compromised agent
simultaneously, so one optimisation step trades the agents off against one another.

\subsubsection*{Differentiating through the routing decision}

The environment decodes a routing decision by taking, for each destination $d$, the
$\arg\max$ over the $K$ candidate paths. The realised route is therefore a piecewise
constant function of the observation, and its gradient vanishes almost everywhere,
which would leave Equation~\eqref{eq:joint_attack} without usable first-order
information. We adopt a straight-through estimator: the forward pass uses the hard
per-destination one-hot decision actually executed by the environment, while the
backward pass propagates the gradient of the underlying softmax policy,
\begin{equation}
  \mathrm{ST}\big(\pi_{\theta_i}(\hat{s})\big)
  = \operatorname{sg}\!\Big[
      \operatorname{onehot}\big(\pi_{\theta_i}(\hat{s})\big) - \pi_{\theta_i}(\hat{s})
    \Big] + \pi_{\theta_i}(\hat{s}),
  \label{eq:straight_through}
\end{equation}

\noindent where $\operatorname{sg}[\cdot]$ denotes the stop-gradient operator and
$\operatorname{onehot}(\cdot)$ is applied blockwise, one block of $K$ entries per
destination. This choice matters for interpreting the results: the gradient is
informative about how to \emph{change a routing decision}, which is not the same
quantity as how to \emph{cause congestion}, and the two coincide only when the
network cannot absorb the redirected flow.

\subsubsection*{Attack objectives}

Two instantiations of $\mathcal{J}$ are considered. The first is grounded in the
victim's own centralised critics, evaluated on the joint action induced by the
perturbed observations,
\begin{equation}
  \mathcal{J}_{\text{critic}}(\hat{\mathbf{s}})
  = -\frac{1}{N} \sum_{i \in \mathcal{C}}
    Q_{\phi^i}\Big( x_t, \; \big(\tilde{a}_1, \dots, \tilde{a}_N \big) \Big),
  \qquad
  \tilde{a}_j = \mathrm{ST}\big(\pi_{\theta_j}(\hat{s}_j)\big),
  \label{eq:obj_critic}
\end{equation}

\noindent where $x_t$ is the true centralised state, held fixed during the attack.
Because the centralised critic of the CTDE formulation
(Section~\ref{subsec:actor_critic}) takes the actions of \emph{all} agents as
input~\cite{lowe2017multi}, it is by construction a non-separable function of the
joint perturbation, and ascending $-Q$ steers the joint action towards the region
the victim's own value model rates as worst.

The second objective targets the shared-bottleneck mechanism explicitly. Let
$u^{(i)}_{d,k}$ denote the true bottleneck utilisation of the $k$-th candidate path
to destination $d$ for agent $i$, read from the \emph{clean} observation, and let
$[\tilde{a}_i]_{d,k}$ be the corresponding entry of the straight-through action.
Then
\begin{equation}
  \mathcal{J}_{\text{cong}}(\hat{\mathbf{s}})
  = \sum_{i \in \mathcal{C}} \sum_{d} \sum_{k}
    \big[\tilde{a}_i\big]_{d,k} \; u^{(i)}_{d,k}.
  \label{eq:obj_congestion}
\end{equation}

\noindent Summing over agents rewards configurations in which several agents are
steered onto the same congested path, which is the mechanism under test. The
utilisation coefficients are taken from the clean observation rather than the
perturbed one; were they read from $\hat{\mathbf{s}}$, the attacker would be
optimising against its own forged view of the network instead of its true state.

\subsubsection*{Joint projected gradient ascent}

Equation~\eqref{eq:joint_attack} is solved by projected gradient ascent following
the standard PGD construction~\cite{madry2019deeplearningmodelsresistant}, applied
to the stacked perturbation. Starting from a uniformly random point inside the ball,
\begin{align}
  \hat{\mathbf{s}}^{(0)}   & = \Pi_{\mathcal{S}}\big( \mathbf{s}_t + \boldsymbol{\delta}_0 \big),
    \qquad \boldsymbol{\delta}_0 \sim \mathrm{Unif}\big(\mathcal{S}(0,\varepsilon)^{N}\big), \\
  \hat{\mathbf{s}}^{(\tau+1)} & = \Pi_{\mathcal{S}}\Big(
      \hat{\mathbf{s}}^{(\tau)} + \alpha \cdot
      \operatorname{sign}\big( \nabla_{\hat{\mathbf{s}}} \mathcal{J}(\hat{\mathbf{s}}^{(\tau)}) \big)
    \Big).
  \label{eq:joint_pgd_iteration}
\end{align}

The random initialisation is what distinguishes PGD from the iterative FGSM variant,
which starts from the clean point~\cite{kurakin2017adversarialmachinelearningscale};
the procedure is repeated from $R$ independent random starts and the iterate
achieving the highest objective is retained, since the local maxima of the inner
problem are distinct but concentrated in value~\cite{madry2019deeplearningmodelsresistant}.
The step size defaults to $\alpha = 2.5\,\varepsilon / T$ for $T$ iterations, which
guarantees the boundary of the ball is reachable from any interior starting
point~\cite{madry2019deeplearningmodelsresistant}. The full procedure is given in
Algorithm~\ref{alg:coordinated_pgd}.

\begin{algorithm}[H]
\SetAlgoLined
\DontPrintSemicolon
\KwIn{clean observations $\mathbf{s}_t$; centralised state $x_t$; budget $\varepsilon$;
      steps $T$; step size $\alpha$; restarts $R$; objective $\mathcal{J}$}
\KwOut{adversarial observations $\hat{s}_{i,t}$ for all $i \in \mathcal{C}$}
$\mathcal{J}^{\text{best}} \leftarrow -\infty$\;
$U \leftarrow$ path utilisations read from the clean $\mathbf{s}_t$\;
\For{$r \leftarrow 1$ \KwTo $R$}{
  $\hat{\mathbf{s}} \leftarrow \Pi_{\mathcal{S}}\big(\mathbf{s}_t + \boldsymbol{\delta}_0\big)$,
    \quad $\boldsymbol{\delta}_0 \sim \mathrm{Unif}(\mathcal{S}(0,\varepsilon)^N)$
    \tcp*{random start}
  \For{$\tau \leftarrow 1$ \KwTo $T$}{
    $\tilde{a}_j \leftarrow \mathrm{ST}\big(\pi_{\theta_j}(\hat{s}_j)\big)$
      for all $j \in \mathcal{C}$ \tcp*{one shared graph}
    $J \leftarrow \mathcal{J}\big(\tilde{a}_{1:N}, x_t, U\big)$\;
    $g \leftarrow \nabla_{\hat{\mathbf{s}}} J$
      \tcp*{one backward $\Rightarrow$ gradient for every agent}
    $\hat{\mathbf{s}} \leftarrow \Pi_{\mathcal{S}}\big(\hat{\mathbf{s}} + \alpha \operatorname{sign}(g)\big)$\;
  }
  \If{$\mathcal{J}(\hat{\mathbf{s}}) > \mathcal{J}^{\text{best}}$}{
    $\mathcal{J}^{\text{best}} \leftarrow \mathcal{J}(\hat{\mathbf{s}})$;\quad
    $\hat{\mathbf{s}}^{\star} \leftarrow \hat{\mathbf{s}}$\;
  }
}
\Return rows of $\hat{\mathbf{s}}^{\star}$\;
\caption{Coordinated multi-agent PGD (one time step)}
\label{alg:coordinated_pgd}
\end{algorithm}

The attack is solved once per time step for the whole compromised group; each agent
subsequently receives its own block of the joint solution, so the $N$ agents act on
a single coordinated decision rather than on $N$ independent ones.

% ─────────────────────────────────────────────────────────────────────────────
\subsection{(B) Strategically-timed attacks under an \texorpdfstring{$L_0$}{L0} budget}
\label{subsec:timed_attack}

The attacks considered so far act at every time step. A more realistic adversary is
constrained in \emph{how often} it can act, which converts the attack into a
resource-allocation problem over the trajectory. We impose an $L_0$ budget: for an
episode of $T_{\text{ep}}$ steps and a budget fraction $b \in (0,1]$,
\begin{equation}
  \frac{1}{T_{\text{ep}}} \sum_{t=1}^{T_{\text{ep}}} \mathbb{1}\big[\text{attack at } t\big]
  \;\le\; b .
  \label{eq:l0_budget}
\end{equation}

Spending this budget uniformly at random would waste it on steps at which the
network has ample spare capacity. We instead spend it at moments of high leverage,
identified by a \emph{critical-state score}. We use the maximum link utilisation
across the topology,
\begin{equation}
  c_t = \max_{(u,v) \in E} \rho_t(u,v),
  \label{eq:critical_score}
\end{equation}

\noindent where $\rho_t(u,v)$ is the utilisation of link $(u,v)$ at time $t$. This
single quantity responds to both regimes of interest: congestion onset raises it
directly, and the loss of a link raises it indirectly by concentrating traffic onto
the surviving paths, so the score is elevated exactly when redundancy is thinnest.

Because the distribution of $c_t$ is not known in advance and drifts within an
episode, the trigger threshold is calibrated online. Let $W_t$ be a sliding window
of the most recent scores. The adversary acts at time $t$ if
\begin{equation}
  c_t \;\ge\; \theta_t,
  \qquad
  \theta_t = Q_{1-b}\big( \{ c_{t'} \}_{t' \in W_t} \big),
  \label{eq:timing_threshold}
\end{equation}

\noindent where $Q_{1-b}$ denotes the empirical $(1-b)$-quantile. The threshold is
computed from the window \emph{before} $c_t$ is inserted, so that a step never
influences its own decision, and the quantile construction makes the realised
attack rate track $b$ without requiring the remainder of the episode to be known in
advance. A short warm-up period, during which the adversary always acts, seeds the
window before the quantile is meaningful. The procedure is summarised in
Algorithm~\ref{alg:timing_gate}.

\begin{algorithm}[H]
\SetAlgoLined
\DontPrintSemicolon
\KwIn{budget $b$; window size $W$; warm-up length $H$; network state at time $t$}
\KwOut{decision $\textit{fire} \in \{\text{true}, \text{false}\}$}
$c_t \leftarrow \max_{(u,v) \in E} \rho_t(u,v)$ \tcp*{critical-state score}
\eIf{$|\mathcal{H}| < H$}{
  $\textit{fire} \leftarrow \text{true}$ \tcp*{warm-up: seed the window}
}{
  $\theta_t \leftarrow Q_{1-b}(\mathcal{H})$ \tcp*{computed before inserting $c_t$}
  $\textit{fire} \leftarrow (c_t \ge \theta_t)$\;
}
insert $c_t$ into $\mathcal{H}$ (ring buffer of size $W$)\;
\Return \textit{fire}\;
\caption{Critical-state timing gate}
\label{alg:timing_gate}
\end{algorithm}

The gate composes with either attack scope: it decides \emph{whether} to act at a
given step, while Section~\ref{subsec:coordinated_attack} decides \emph{how} the
perturbation is distributed across agents once the decision to act has been taken.
On skipped steps the agents receive their unmodified observations.

% ─────────────────────────────────────────────────────────────────────────────
\subsection{Evaluation protocol}
\label{subsec:eval_protocol}

\subsubsection*{Paired arms}

Each configuration is evaluated with three arms: a \emph{clean} arm with no
perturbation, a \emph{random} control that draws a uniformly random direction from
the same admissible set of Equation~\eqref{eq:admissible_set}, and the \emph{attack}
under test. All three are run on identical episodes: the traffic sequence and any
injected link failures are generated from a common seed, so episode $e$ of one arm
and episode $e$ of another differ only in the perturbation applied. The random
control draws from a dedicated random stream so that it consumes no entropy from
the traffic generator, which would otherwise desynchronise the arms and invalidate
the pairing.

\subsubsection*{Adversarial-specific damage}

Under link failures the raw degradation relative to the clean arm is dominated by
the failures themselves rather than by the attack. The primary reported quantity is
therefore the \emph{adversarial gap}, the paired per-episode difference between the
random control and the attack,
\begin{equation}
  \Delta^{\text{adv}} = \frac{1}{n} \sum_{e=1}^{n}
    \Big( \mathrm{PDR}^{\text{rand}}_{e} - \mathrm{PDR}^{\text{atk}}_{e} \Big),
  \label{eq:adversarial_gap}
\end{equation}

\noindent reported with a paired $95\%$ confidence interval,
$\Delta^{\text{adv}} \pm t_{0.975,\,n-1}\, s_d / \sqrt{n}$, where $s_d$ is the
standard deviation of the per-episode differences. Differencing within episode pairs
removes the episode-to-episode variance that otherwise dominates the comparison. The
raw drop with respect to the clean arm is reported alongside it, but is not used as
the measure of attack strength.

\subsubsection*{Decisions versus outcomes}

Two families of metric are recorded and reported separately. \emph{Decision} metrics
count how much routing behaviour the attack changed, principally the action flip
rate, the fraction of per-agent, per-destination $\arg\max$ path choices that differ
between the clean and adversarial observations at the same time step,
\begin{equation}
  \mathrm{flip}
  = \frac{1}{|\mathcal{C}|\, n_d\, T_{\text{ep}}}
    \sum_{t}\sum_{i \in \mathcal{C}}\sum_{d}
    \mathbb{1}\Big[
      \arg\max_k \pi_{\theta_i}\big(\hat{s}_{i,t}\big)_{d,k}
      \neq
      \arg\max_k \pi_{\theta_i}\big(s_{i,t}^{\text{local}}\big)_{d,k}
    \Big].
  \label{eq:flip_rate}
\end{equation}

\emph{Outcome} metrics measure delivered traffic, principally end-to-end PDR. The
distinction is essential in a topology with $K$-path redundancy, where a high flip
rate accompanied by an unchanged PDR indicates that the network absorbed the
redirected flows rather than that the attack succeeded. Reporting a flip rate as
though it were a measure of damage would conflate the two.
.
%%%%%%%%%%%%%%%%%%%%%%%%%%%%%%%%%%%%%%%%%%%%%%%%%%%%%%%%%%%%%%%%%%%%

% ─────────────────────────────────────────────────────────────────────────────
\subsection{Threat model and admissible perturbation set}
\label{subsec:threat_model_pgd}

Both extensions share a single threat model, which is deliberately identical to the
one used for the FGSM baseline of Section~\ref{subsec:fgsm_maddpg_def} so that any
observed difference in damage is attributable to the attack itself rather than to a
change in the attacker's capabilities.

Let $\mathcal{C} \subseteq \{1,\dots,M\}$ denote the set of compromised agents and
$N = |\mathcal{C}|$ its cardinality. The victim policies $\{\pi_{\theta_i}\}$ and
critics $\{Q_{\phi^i}\}$ are \emph{frozen}: their parameters are fixed throughout,
and the adversary interacts with them only through forward evaluations and
gradients taken with respect to the \emph{input} observations. No victim parameter
is ever updated, so the attack measures the robustness of a fixed trained policy
rather than the outcome of an adversarial training game.

The adversary perturbs observations only; the underlying network state evolves from
the true, unperturbed dynamics. For agent $i$ at time $t$ the admissible set is
\begin{equation}
  \mathcal{S}\big(s_{i,t}^{\text{local}}, \varepsilon\big) =
  \Big\{ \hat{s} \;\Big|\;
    \big\| \hat{s} - s_{i,t}^{\text{local}} \big\|_{\infty} \le \varepsilon,
    \;\; \hat{s} \in \mathcal{D}
  \Big\},
  \label{eq:admissible_set}
\end{equation}

\noindent where $\mathcal{D}$ is the domain constraint that re-projects the
normalised bandwidth features onto $[0,1]$. The corresponding projection operator
$\Pi_{\mathcal{S}}$ therefore composes an $\ell_{\infty}$ clip about the clean
observation with the domain clip of Equation~\eqref{eq:fgsm_maddpg_obs}. Because
the measured damage depends on the size of the admissible set, the \emph{same}
operator $\Pi_{\mathcal{S}}$ is applied to every attack and to the random control
of Section~\ref{subsec:eval_protocol}; using a stricter domain clip for one attack
than another would confound the perturbation direction with the size of the set it
is drawn from.

% ─────────────────────────────────────────────────────────────────────────────
\subsection{(A) Coordinated multi-agent perturbation}
\label{subsec:coordinated_attack}

\subsubsection*{Separability of the per-agent attack}

The FGSM baseline, and any attack that iterates over agents independently, solves
$N$ decoupled problems: for each $i \in \mathcal{C}$,
\begin{equation}
  \eta_{i,t}^{\star} = \arg\max_{\eta \,:\, s_{i,t}^{\text{local}} + \eta \in \mathcal{S}}
  \; \mathcal{L}_{i}\big(s_{i,t}^{\text{local}} + \eta\big),
  \label{eq:separable_attack}
\end{equation}

\noindent where $\mathcal{L}_i$ depends only on agent $i$'s own observation and
policy. The joint perturbation obtained this way is the concatenation of $N$
independent solutions, and the implied joint objective is \emph{separable},
\begin{equation}
  \mathcal{J}_{\text{sep}}(\hat{\mathbf{s}}_t)
  = \sum_{i \in \mathcal{C}} \mathcal{L}_i(\hat{s}_{i,t}),
  \label{eq:separable_objective}
\end{equation}

\noindent so no term couples two agents. This is a structural limitation rather
than a matter of attack strength: a separable objective cannot express a preference
for outcomes that depend on \emph{which} agents are steered onto the \emph{same}
resource, because the value of steering agent $i$ onto a link is independent of
whether agent $j$ was steered onto it as well. In a topology whose robustness comes
from $K$-path redundancy, the damaging configurations are precisely those in which
several redirected flows converge on one surviving link, and these lie outside the
reach of Equation~\eqref{eq:separable_objective}.

\subsubsection*{Joint formulation}

We therefore optimise over the joint perturbation directly. Writing
$\mathbf{s}_t = \big(s_{1,t}^{\text{local}}, \dots, s_{N,t}^{\text{local}}\big)$ for
the stacked observations of the compromised agents, the coordinated attack solves
\begin{equation}
  \hat{\mathbf{s}}_t^{\star}
  = \arg\max_{\hat{\mathbf{s}} \,\in\, \mathcal{S}(\mathbf{s}_t, \varepsilon)^{N}}
  \; \mathcal{J}\big(\hat{\mathbf{s}}\big),
  \label{eq:joint_attack}
\end{equation}

\noindent with a \emph{non-separable} global objective $\mathcal{J}$. A single
gradient of $\mathcal{J}$ supplies a descent direction for every compromised agent
simultaneously, so one optimisation step trades the agents off against one another.

\subsubsection*{Differentiating through the routing decision}

The environment decodes a routing decision by taking, for each destination $d$, the
$\arg\max$ over the $K$ candidate paths. The realised route is therefore a piecewise
constant function of the observation, and its gradient vanishes almost everywhere,
which would leave Equation~\eqref{eq:joint_attack} without usable first-order
information. We adopt a straight-through estimator: the forward pass uses the hard
per-destination one-hot decision actually executed by the environment, while the
backward pass propagates the gradient of the underlying softmax policy,
\begin{equation}
  \mathrm{ST}\big(\pi_{\theta_i}(\hat{s})\big)
  = \operatorname{sg}\!\Big[
      \operatorname{onehot}\big(\pi_{\theta_i}(\hat{s})\big) - \pi_{\theta_i}(\hat{s})
    \Big] + \pi_{\theta_i}(\hat{s}),
  \label{eq:straight_through}
\end{equation}

\noindent where $\operatorname{sg}[\cdot]$ denotes the stop-gradient operator and
$\operatorname{onehot}(\cdot)$ is applied blockwise, one block of $K$ entries per
destination. This choice matters for interpreting the results: the gradient is
informative about how to \emph{change a routing decision}, which is not the same
quantity as how to \emph{cause congestion}, and the two coincide only when the
network cannot absorb the redirected flow.

\subsubsection*{Attack objectives}

Two instantiations of $\mathcal{J}$ are considered. The first is grounded in the
victim's own centralised critics, evaluated on the joint action induced by the
perturbed observations,
\begin{equation}
  \mathcal{J}_{\text{critic}}(\hat{\mathbf{s}})
  = -\frac{1}{N} \sum_{i \in \mathcal{C}}
    Q_{\phi^i}\Big( x_t, \; \big(\tilde{a}_1, \dots, \tilde{a}_N \big) \Big),
  \qquad
  \tilde{a}_j = \mathrm{ST}\big(\pi_{\theta_j}(\hat{s}_j)\big),
  \label{eq:obj_critic}
\end{equation}

\noindent where $x_t$ is the true centralised state, held fixed during the attack.
Because the centralised critic of the CTDE formulation
(Section~\ref{subsec:actor_critic}) takes the actions of \emph{all} agents as
input~\cite{lowe2017multi}, it is by construction a non-separable function of the
joint perturbation, and ascending $-Q$ steers the joint action towards the region
the victim's own value model rates as worst.

The second objective targets the shared-bottleneck mechanism explicitly. Let
$u^{(i)}_{d,k}$ denote the true bottleneck utilisation of the $k$-th candidate path
to destination $d$ for agent $i$, read from the \emph{clean} observation, and let
$[\tilde{a}_i]_{d,k}$ be the corresponding entry of the straight-through action.
Then
\begin{equation}
  \mathcal{J}_{\text{cong}}(\hat{\mathbf{s}})
  = \sum_{i \in \mathcal{C}} \sum_{d} \sum_{k}
    \big[\tilde{a}_i\big]_{d,k} \; u^{(i)}_{d,k}.
  \label{eq:obj_congestion}
\end{equation}

\noindent Summing over agents rewards configurations in which several agents are
steered onto the same congested path, which is the mechanism under test. The
utilisation coefficients are taken from the clean observation rather than the
perturbed one; were they read from $\hat{\mathbf{s}}$, the attacker would be
optimising against its own forged view of the network instead of its true state.

\subsubsection*{Joint projected gradient ascent}

Equation~\eqref{eq:joint_attack} is solved by projected gradient ascent following
the standard PGD construction~\cite{madry2019deeplearningmodelsresistant}, applied
to the stacked perturbation. Starting from a uniformly random point inside the ball,
\begin{align}
  \hat{\mathbf{s}}^{(0)}   & = \Pi_{\mathcal{S}}\big( \mathbf{s}_t + \boldsymbol{\delta}_0 \big),
    \qquad \boldsymbol{\delta}_0 \sim \mathrm{Unif}\big(\mathcal{S}(0,\varepsilon)^{N}\big), \\
  \hat{\mathbf{s}}^{(\tau+1)} & = \Pi_{\mathcal{S}}\Big(
      \hat{\mathbf{s}}^{(\tau)} + \alpha \cdot
      \operatorname{sign}\big( \nabla_{\hat{\mathbf{s}}} \mathcal{J}(\hat{\mathbf{s}}^{(\tau)}) \big)
    \Big).
  \label{eq:joint_pgd_iteration}
\end{align}

The random initialisation is what distinguishes PGD from the iterative FGSM variant,
which starts from the clean point~\cite{kurakin2017adversarialmachinelearningscale};
the procedure is repeated from $R$ independent random starts and the iterate
achieving the highest objective is retained, since the local maxima of the inner
problem are distinct but concentrated in value~\cite{madry2019deeplearningmodelsresistant}.
The step size defaults to $\alpha = 2.5\,\varepsilon / T$ for $T$ iterations, which
guarantees the boundary of the ball is reachable from any interior starting
point~\cite{madry2019deeplearningmodelsresistant}. The full procedure is given in
Algorithm~\ref{alg:coordinated_pgd}.

\begin{algorithm}[H]
\SetAlgoLined
\DontPrintSemicolon
\KwIn{clean observations $\mathbf{s}_t$; centralised state $x_t$; budget $\varepsilon$;
      steps $T$; step size $\alpha$; restarts $R$; objective $\mathcal{J}$}
\KwOut{adversarial observations $\hat{s}_{i,t}$ for all $i \in \mathcal{C}$}
$\mathcal{J}^{\text{best}} \leftarrow -\infty$\;
$U \leftarrow$ path utilisations read from the clean $\mathbf{s}_t$\;
\For{$r \leftarrow 1$ \KwTo $R$}{
  $\hat{\mathbf{s}} \leftarrow \Pi_{\mathcal{S}}\big(\mathbf{s}_t + \boldsymbol{\delta}_0\big)$,
    \quad $\boldsymbol{\delta}_0 \sim \mathrm{Unif}(\mathcal{S}(0,\varepsilon)^N)$
    \tcp*{random start}
  \For{$\tau \leftarrow 1$ \KwTo $T$}{
    $\tilde{a}_j \leftarrow \mathrm{ST}\big(\pi_{\theta_j}(\hat{s}_j)\big)$
      for all $j \in \mathcal{C}$ \tcp*{one shared graph}
    $J \leftarrow \mathcal{J}\big(\tilde{a}_{1:N}, x_t, U\big)$\;
    $g \leftarrow \nabla_{\hat{\mathbf{s}}} J$
      \tcp*{one backward $\Rightarrow$ gradient for every agent}
    $\hat{\mathbf{s}} \leftarrow \Pi_{\mathcal{S}}\big(\hat{\mathbf{s}} + \alpha \operatorname{sign}(g)\big)$\;
  }
  \If{$\mathcal{J}(\hat{\mathbf{s}}) > \mathcal{J}^{\text{best}}$}{
    $\mathcal{J}^{\text{best}} \leftarrow \mathcal{J}(\hat{\mathbf{s}})$;\quad
    $\hat{\mathbf{s}}^{\star} \leftarrow \hat{\mathbf{s}}$\;
  }
}
\Return rows of $\hat{\mathbf{s}}^{\star}$\;
\caption{Coordinated multi-agent PGD (one time step)}
\label{alg:coordinated_pgd}
\end{algorithm}

The attack is solved once per time step for the whole compromised group; each agent
subsequently receives its own block of the joint solution, so the $N$ agents act on
a single coordinated decision rather than on $N$ independent ones.

% ─────────────────────────────────────────────────────────────────────────────
\subsection{(B) Strategically-timed attacks under an \texorpdfstring{$L_0$}{L0} budget}
\label{subsec:timed_attack}

The attacks considered so far act at every time step. A more realistic adversary is
constrained in \emph{how often} it can act, which converts the attack into a
resource-allocation problem over the trajectory. We impose an $L_0$ budget: for an
episode of $T_{\text{ep}}$ steps and a budget fraction $b \in (0,1]$,
\begin{equation}
  \frac{1}{T_{\text{ep}}} \sum_{t=1}^{T_{\text{ep}}} \mathbb{1}\big[\text{attack at } t\big]
  \;\le\; b .
  \label{eq:l0_budget}
\end{equation}

Spending this budget uniformly at random would waste it on steps at which the
network has ample spare capacity. We instead spend it at moments of high leverage,
identified by a \emph{critical-state score}. We use the maximum link utilisation
across the topology,
\begin{equation}
  c_t = \max_{(u,v) \in E} \rho_t(u,v),
  \label{eq:critical_score}
\end{equation}

\noindent where $\rho_t(u,v)$ is the utilisation of link $(u,v)$ at time $t$. This
single quantity responds to both regimes of interest: congestion onset raises it
directly, and the loss of a link raises it indirectly by concentrating traffic onto
the surviving paths, so the score is elevated exactly when redundancy is thinnest.

Because the distribution of $c_t$ is not known in advance and drifts within an
episode, the trigger threshold is calibrated online. Let $W_t$ be a sliding window
of the most recent scores. The adversary acts at time $t$ if
\begin{equation}
  c_t \;\ge\; \theta_t,
  \qquad
  \theta_t = Q_{1-b}\big( \{ c_{t'} \}_{t' \in W_t} \big),
  \label{eq:timing_threshold}
\end{equation}

\noindent where $Q_{1-b}$ denotes the empirical $(1-b)$-quantile. The threshold is
computed from the window \emph{before} $c_t$ is inserted, so that a step never
influences its own decision, and the quantile construction makes the realised
attack rate track $b$ without requiring the remainder of the episode to be known in
advance. A short warm-up period, during which the adversary always acts, seeds the
window before the quantile is meaningful. The procedure is summarised in
Algorithm~\ref{alg:timing_gate}.

\begin{algorithm}[H]
\SetAlgoLined
\DontPrintSemicolon
\KwIn{budget $b$; window size $W$; warm-up length $H$; network state at time $t$}
\KwOut{decision $\textit{fire} \in \{\text{true}, \text{false}\}$}
$c_t \leftarrow \max_{(u,v) \in E} \rho_t(u,v)$ \tcp*{critical-state score}
\eIf{$|\mathcal{H}| < H$}{
  $\textit{fire} \leftarrow \text{true}$ \tcp*{warm-up: seed the window}
}{
  $\theta_t \leftarrow Q_{1-b}(\mathcal{H})$ \tcp*{computed before inserting $c_t$}
  $\textit{fire} \leftarrow (c_t \ge \theta_t)$\;
}
insert $c_t$ into $\mathcal{H}$ (ring buffer of size $W$)\;
\Return \textit{fire}\;
\caption{Critical-state timing gate}
\label{alg:timing_gate}
\end{algorithm}

The gate composes with either attack scope: it decides \emph{whether} to act at a
given step, while Section~\ref{subsec:coordinated_attack} decides \emph{how} the
perturbation is distributed across agents once the decision to act has been taken.
On skipped steps the agents receive their unmodified observations.

% ─────────────────────────────────────────────────────────────────────────────
\subsection{Evaluation protocol}
\label{subsec:eval_protocol}

\subsubsection*{Paired arms}

Each configuration is evaluated with three arms: a \emph{clean} arm with no
perturbation, a \emph{random} control that draws a uniformly random direction from
the same admissible set of Equation~\eqref{eq:admissible_set}, and the \emph{attack}
under test. All three are run on identical episodes: the traffic sequence and any
injected link failures are generated from a common seed, so episode $e$ of one arm
and episode $e$ of another differ only in the perturbation applied. The random
control draws from a dedicated random stream so that it consumes no entropy from
the traffic generator, which would otherwise desynchronise the arms and invalidate
the pairing.

\subsubsection*{Adversarial-specific damage}

Under link failures the raw degradation relative to the clean arm is dominated by
the failures themselves rather than by the attack. The primary reported quantity is
therefore the \emph{adversarial gap}, the paired per-episode difference between the
random control and the attack,
\begin{equation}
  \Delta^{\text{adv}} = \frac{1}{n} \sum_{e=1}^{n}
    \Big( \mathrm{PDR}^{\text{rand}}_{e} - \mathrm{PDR}^{\text{atk}}_{e} \Big),
  \label{eq:adversarial_gap}
\end{equation}

\noindent reported with a paired $95\%$ confidence interval,
$\Delta^{\text{adv}} \pm t_{0.975,\,n-1}\, s_d / \sqrt{n}$, where $s_d$ is the
standard deviation of the per-episode differences. Differencing within episode pairs
removes the episode-to-episode variance that otherwise dominates the comparison. The
raw drop with respect to the clean arm is reported alongside it, but is not used as
the measure of attack strength.

\subsubsection*{Decisions versus outcomes}

Two families of metric are recorded and reported separately. \emph{Decision} metrics
count how much routing behaviour the attack changed, principally the action flip
rate, the fraction of per-agent, per-destination $\arg\max$ path choices that differ
between the clean and adversarial observations at the same time step,
\begin{equation}
  \mathrm{flip}
  = \frac{1}{|\mathcal{C}|\, n_d\, T_{\text{ep}}}
    \sum_{t}\sum_{i \in \mathcal{C}}\sum_{d}
    \mathbb{1}\Big[
      \arg\max_k \pi_{\theta_i}\big(\hat{s}_{i,t}\big)_{d,k}
      \neq
      \arg\max_k \pi_{\theta_i}\big(s_{i,t}^{\text{local}}\big)_{d,k}
    \Big].
  \label{eq:flip_rate}
\end{equation}

\emph{Outcome} metrics measure delivered traffic, principally end-to-end PDR. The
distinction is essential in a topology with $K$-path redundancy, where a high flip
rate accompanied by an unchanged PDR indicates that the network absorbed the
redirected flows rather than that the attack succeeded. Reporting a flip rate as
though it were a measure of damage would conflate the two.
.
%%%%%%%%%%%%%%%%%%%%%%%%%%%%%%%%%%%%%%%%%%%%%%%%%%%%%%%%%%%%%%%%%%%%

% ─────────────────────────────────────────────────────────────────────────────
\subsection{Threat model and admissible perturbation set}
\label{subsec:threat_model_pgd}

Both extensions share a single threat model, which is deliberately identical to the
one used for the FGSM baseline of Section~\ref{subsec:fgsm_maddpg_def} so that any
observed difference in damage is attributable to the attack itself rather than to a
change in the attacker's capabilities.

Let $\mathcal{C} \subseteq \{1,\dots,M\}$ denote the set of compromised agents and
$N = |\mathcal{C}|$ its cardinality. The victim policies $\{\pi_{\theta_i}\}$ and
critics $\{Q_{\phi^i}\}$ are \emph{frozen}: their parameters are fixed throughout,
and the adversary interacts with them only through forward evaluations and
gradients taken with respect to the \emph{input} observations. No victim parameter
is ever updated, so the attack measures the robustness of a fixed trained policy
rather than the outcome of an adversarial training game.

The adversary perturbs observations only; the underlying network state evolves from
the true, unperturbed dynamics. For agent $i$ at time $t$ the admissible set is
\begin{equation}
  \mathcal{S}\big(s_{i,t}^{\text{local}}, \varepsilon\big) =
  \Big\{ \hat{s} \;\Big|\;
    \big\| \hat{s} - s_{i,t}^{\text{local}} \big\|_{\infty} \le \varepsilon,
    \;\; \hat{s} \in \mathcal{D}
  \Big\},
  \label{eq:admissible_set}
\end{equation}

\noindent where $\mathcal{D}$ is the domain constraint that re-projects the
normalised bandwidth features onto $[0,1]$. The corresponding projection operator
$\Pi_{\mathcal{S}}$ therefore composes an $\ell_{\infty}$ clip about the clean
observation with the domain clip of Equation~\eqref{eq:fgsm_maddpg_obs}. Because
the measured damage depends on the size of the admissible set, the \emph{same}
operator $\Pi_{\mathcal{S}}$ is applied to every attack and to the random control
of Section~\ref{subsec:eval_protocol}; using a stricter domain clip for one attack
than another would confound the perturbation direction with the size of the set it
is drawn from.

% ─────────────────────────────────────────────────────────────────────────────
\subsection{(A) Coordinated multi-agent perturbation}
\label{subsec:coordinated_attack}

\subsubsection*{Separability of the per-agent attack}

The FGSM baseline, and any attack that iterates over agents independently, solves
$N$ decoupled problems: for each $i \in \mathcal{C}$,
\begin{equation}
  \eta_{i,t}^{\star} = \arg\max_{\eta \,:\, s_{i,t}^{\text{local}} + \eta \in \mathcal{S}}
  \; \mathcal{L}_{i}\big(s_{i,t}^{\text{local}} + \eta\big),
  \label{eq:separable_attack}
\end{equation}

\noindent where $\mathcal{L}_i$ depends only on agent $i$'s own observation and
policy. The joint perturbation obtained this way is the concatenation of $N$
independent solutions, and the implied joint objective is \emph{separable},
\begin{equation}
  \mathcal{J}_{\text{sep}}(\hat{\mathbf{s}}_t)
  = \sum_{i \in \mathcal{C}} \mathcal{L}_i(\hat{s}_{i,t}),
  \label{eq:separable_objective}
\end{equation}

\noindent so no term couples two agents. This is a structural limitation rather
than a matter of attack strength: a separable objective cannot express a preference
for outcomes that depend on \emph{which} agents are steered onto the \emph{same}
resource, because the value of steering agent $i$ onto a link is independent of
whether agent $j$ was steered onto it as well. In a topology whose robustness comes
from $K$-path redundancy, the damaging configurations are precisely those in which
several redirected flows converge on one surviving link, and these lie outside the
reach of Equation~\eqref{eq:separable_objective}.

\subsubsection*{Joint formulation}

We therefore optimise over the joint perturbation directly. Writing
$\mathbf{s}_t = \big(s_{1,t}^{\text{local}}, \dots, s_{N,t}^{\text{local}}\big)$ for
the stacked observations of the compromised agents, the coordinated attack solves
\begin{equation}
  \hat{\mathbf{s}}_t^{\star}
  = \arg\max_{\hat{\mathbf{s}} \,\in\, \mathcal{S}(\mathbf{s}_t, \varepsilon)^{N}}
  \; \mathcal{J}\big(\hat{\mathbf{s}}\big),
  \label{eq:joint_attack}
\end{equation}

\noindent with a \emph{non-separable} global objective $\mathcal{J}$. A single
gradient of $\mathcal{J}$ supplies a descent direction for every compromised agent
simultaneously, so one optimisation step trades the agents off against one another.

\subsubsection*{Differentiating through the routing decision}

The environment decodes a routing decision by taking, for each destination $d$, the
$\arg\max$ over the $K$ candidate paths. The realised route is therefore a piecewise
constant function of the observation, and its gradient vanishes almost everywhere,
which would leave Equation~\eqref{eq:joint_attack} without usable first-order
information. We adopt a straight-through estimator: the forward pass uses the hard
per-destination one-hot decision actually executed by the environment, while the
backward pass propagates the gradient of the underlying softmax policy,
\begin{equation}
  \mathrm{ST}\big(\pi_{\theta_i}(\hat{s})\big)
  = \operatorname{sg}\!\Big[
      \operatorname{onehot}\big(\pi_{\theta_i}(\hat{s})\big) - \pi_{\theta_i}(\hat{s})
    \Big] + \pi_{\theta_i}(\hat{s}),
  \label{eq:straight_through}
\end{equation}

\noindent where $\operatorname{sg}[\cdot]$ denotes the stop-gradient operator and
$\operatorname{onehot}(\cdot)$ is applied blockwise, one block of $K$ entries per
destination. This choice matters for interpreting the results: the gradient is
informative about how to \emph{change a routing decision}, which is not the same
quantity as how to \emph{cause congestion}, and the two coincide only when the
network cannot absorb the redirected flow.

\subsubsection*{Attack objectives}

Two instantiations of $\mathcal{J}$ are considered. The first is grounded in the
victim's own centralised critics, evaluated on the joint action induced by the
perturbed observations,
\begin{equation}
  \mathcal{J}_{\text{critic}}(\hat{\mathbf{s}})
  = -\frac{1}{N} \sum_{i \in \mathcal{C}}
    Q_{\phi^i}\Big( x_t, \; \big(\tilde{a}_1, \dots, \tilde{a}_N \big) \Big),
  \qquad
  \tilde{a}_j = \mathrm{ST}\big(\pi_{\theta_j}(\hat{s}_j)\big),
  \label{eq:obj_critic}
\end{equation}

\noindent where $x_t$ is the true centralised state, held fixed during the attack.
Because the centralised critic of the CTDE formulation
(Section~\ref{subsec:actor_critic}) takes the actions of \emph{all} agents as
input~\cite{lowe2017multi}, it is by construction a non-separable function of the
joint perturbation, and ascending $-Q$ steers the joint action towards the region
the victim's own value model rates as worst.

The second objective targets the shared-bottleneck mechanism explicitly. Let
$u^{(i)}_{d,k}$ denote the true bottleneck utilisation of the $k$-th candidate path
to destination $d$ for agent $i$, read from the \emph{clean} observation, and let
$[\tilde{a}_i]_{d,k}$ be the corresponding entry of the straight-through action.
Then
\begin{equation}
  \mathcal{J}_{\text{cong}}(\hat{\mathbf{s}})
  = \sum_{i \in \mathcal{C}} \sum_{d} \sum_{k}
    \big[\tilde{a}_i\big]_{d,k} \; u^{(i)}_{d,k}.
  \label{eq:obj_congestion}
\end{equation}

\noindent Summing over agents rewards configurations in which several agents are
steered onto the same congested path, which is the mechanism under test. The
utilisation coefficients are taken from the clean observation rather than the
perturbed one; were they read from $\hat{\mathbf{s}}$, the attacker would be
optimising against its own forged view of the network instead of its true state.

\subsubsection*{Joint projected gradient ascent}

Equation~\eqref{eq:joint_attack} is solved by projected gradient ascent following
the standard PGD construction~\cite{madry2019deeplearningmodelsresistant}, applied
to the stacked perturbation. Starting from a uniformly random point inside the ball,
\begin{align}
  \hat{\mathbf{s}}^{(0)}   & = \Pi_{\mathcal{S}}\big( \mathbf{s}_t + \boldsymbol{\delta}_0 \big),
    \qquad \boldsymbol{\delta}_0 \sim \mathrm{Unif}\big(\mathcal{S}(0,\varepsilon)^{N}\big), \\
  \hat{\mathbf{s}}^{(\tau+1)} & = \Pi_{\mathcal{S}}\Big(
      \hat{\mathbf{s}}^{(\tau)} + \alpha \cdot
      \operatorname{sign}\big( \nabla_{\hat{\mathbf{s}}} \mathcal{J}(\hat{\mathbf{s}}^{(\tau)}) \big)
    \Big).
  \label{eq:joint_pgd_iteration}
\end{align}

The random initialisation is what distinguishes PGD from the iterative FGSM variant,
which starts from the clean point~\cite{kurakin2017adversarialmachinelearningscale};
the procedure is repeated from $R$ independent random starts and the iterate
achieving the highest objective is retained, since the local maxima of the inner
problem are distinct but concentrated in value~\cite{madry2019deeplearningmodelsresistant}.
The step size defaults to $\alpha = 2.5\,\varepsilon / T$ for $T$ iterations, which
guarantees the boundary of the ball is reachable from any interior starting
point~\cite{madry2019deeplearningmodelsresistant}. The full procedure is given in
Algorithm~\ref{alg:coordinated_pgd}.

\begin{algorithm}[H]
\SetAlgoLined
\DontPrintSemicolon
\KwIn{clean observations $\mathbf{s}_t$; centralised state $x_t$; budget $\varepsilon$;
      steps $T$; step size $\alpha$; restarts $R$; objective $\mathcal{J}$}
\KwOut{adversarial observations $\hat{s}_{i,t}$ for all $i \in \mathcal{C}$}
$\mathcal{J}^{\text{best}} \leftarrow -\infty$\;
$U \leftarrow$ path utilisations read from the clean $\mathbf{s}_t$\;
\For{$r \leftarrow 1$ \KwTo $R$}{
  $\hat{\mathbf{s}} \leftarrow \Pi_{\mathcal{S}}\big(\mathbf{s}_t + \boldsymbol{\delta}_0\big)$,
    \quad $\boldsymbol{\delta}_0 \sim \mathrm{Unif}(\mathcal{S}(0,\varepsilon)^N)$
    \tcp*{random start}
  \For{$\tau \leftarrow 1$ \KwTo $T$}{
    $\tilde{a}_j \leftarrow \mathrm{ST}\big(\pi_{\theta_j}(\hat{s}_j)\big)$
      for all $j \in \mathcal{C}$ \tcp*{one shared graph}
    $J \leftarrow \mathcal{J}\big(\tilde{a}_{1:N}, x_t, U\big)$\;
    $g \leftarrow \nabla_{\hat{\mathbf{s}}} J$
      \tcp*{one backward $\Rightarrow$ gradient for every agent}
    $\hat{\mathbf{s}} \leftarrow \Pi_{\mathcal{S}}\big(\hat{\mathbf{s}} + \alpha \operatorname{sign}(g)\big)$\;
  }
  \If{$\mathcal{J}(\hat{\mathbf{s}}) > \mathcal{J}^{\text{best}}$}{
    $\mathcal{J}^{\text{best}} \leftarrow \mathcal{J}(\hat{\mathbf{s}})$;\quad
    $\hat{\mathbf{s}}^{\star} \leftarrow \hat{\mathbf{s}}$\;
  }
}
\Return rows of $\hat{\mathbf{s}}^{\star}$\;
\caption{Coordinated multi-agent PGD (one time step)}
\label{alg:coordinated_pgd}
\end{algorithm}

The attack is solved once per time step for the whole compromised group; each agent
subsequently receives its own block of the joint solution, so the $N$ agents act on
a single coordinated decision rather than on $N$ independent ones.

% ─────────────────────────────────────────────────────────────────────────────
\subsection{(B) Strategically-timed attacks under an \texorpdfstring{$L_0$}{L0} budget}
\label{subsec:timed_attack}

The attacks considered so far act at every time step. A more realistic adversary is
constrained in \emph{how often} it can act, which converts the attack into a
resource-allocation problem over the trajectory. We impose an $L_0$ budget: for an
episode of $T_{\text{ep}}$ steps and a budget fraction $b \in (0,1]$,
\begin{equation}
  \frac{1}{T_{\text{ep}}} \sum_{t=1}^{T_{\text{ep}}} \mathbb{1}\big[\text{attack at } t\big]
  \;\le\; b .
  \label{eq:l0_budget}
\end{equation}

Spending this budget uniformly at random would waste it on steps at which the
network has ample spare capacity. We instead spend it at moments of high leverage,
identified by a \emph{critical-state score}. We use the maximum link utilisation
across the topology,
\begin{equation}
  c_t = \max_{(u,v) \in E} \rho_t(u,v),
  \label{eq:critical_score}
\end{equation}

\noindent where $\rho_t(u,v)$ is the utilisation of link $(u,v)$ at time $t$. This
single quantity responds to both regimes of interest: congestion onset raises it
directly, and the loss of a link raises it indirectly by concentrating traffic onto
the surviving paths, so the score is elevated exactly when redundancy is thinnest.

Because the distribution of $c_t$ is not known in advance and drifts within an
episode, the trigger threshold is calibrated online. Let $W_t$ be a sliding window
of the most recent scores. The adversary acts at time $t$ if
\begin{equation}
  c_t \;\ge\; \theta_t,
  \qquad
  \theta_t = Q_{1-b}\big( \{ c_{t'} \}_{t' \in W_t} \big),
  \label{eq:timing_threshold}
\end{equation}

\noindent where $Q_{1-b}$ denotes the empirical $(1-b)$-quantile. The threshold is
computed from the window \emph{before} $c_t$ is inserted, so that a step never
influences its own decision, and the quantile construction makes the realised
attack rate track $b$ without requiring the remainder of the episode to be known in
advance. A short warm-up period, during which the adversary always acts, seeds the
window before the quantile is meaningful. The procedure is summarised in
Algorithm~\ref{alg:timing_gate}.

\begin{algorithm}[H]
\SetAlgoLined
\DontPrintSemicolon
\KwIn{budget $b$; window size $W$; warm-up length $H$; network state at time $t$}
\KwOut{decision $\textit{fire} \in \{\text{true}, \text{false}\}$}
$c_t \leftarrow \max_{(u,v) \in E} \rho_t(u,v)$ \tcp*{critical-state score}
\eIf{$|\mathcal{H}| < H$}{
  $\textit{fire} \leftarrow \text{true}$ \tcp*{warm-up: seed the window}
}{
  $\theta_t \leftarrow Q_{1-b}(\mathcal{H})$ \tcp*{computed before inserting $c_t$}
  $\textit{fire} \leftarrow (c_t \ge \theta_t)$\;
}
insert $c_t$ into $\mathcal{H}$ (ring buffer of size $W$)\;
\Return \textit{fire}\;
\caption{Critical-state timing gate}
\label{alg:timing_gate}
\end{algorithm}

The gate composes with either attack scope: it decides \emph{whether} to act at a
given step, while Section~\ref{subsec:coordinated_attack} decides \emph{how} the
perturbation is distributed across agents once the decision to act has been taken.
On skipped steps the agents receive their unmodified observations.

% ─────────────────────────────────────────────────────────────────────────────
\subsection{Evaluation protocol}
\label{subsec:eval_protocol}

\subsubsection*{Paired arms}

Each configuration is evaluated with three arms: a \emph{clean} arm with no
perturbation, a \emph{random} control that draws a uniformly random direction from
the same admissible set of Equation~\eqref{eq:admissible_set}, and the \emph{attack}
under test. All three are run on identical episodes: the traffic sequence and any
injected link failures are generated from a common seed, so episode $e$ of one arm
and episode $e$ of another differ only in the perturbation applied. The random
control draws from a dedicated random stream so that it consumes no entropy from
the traffic generator, which would otherwise desynchronise the arms and invalidate
the pairing.

\subsubsection*{Adversarial-specific damage}

Under link failures the raw degradation relative to the clean arm is dominated by
the failures themselves rather than by the attack. The primary reported quantity is
therefore the \emph{adversarial gap}, the paired per-episode difference between the
random control and the attack,
\begin{equation}
  \Delta^{\text{adv}} = \frac{1}{n} \sum_{e=1}^{n}
    \Big( \mathrm{PDR}^{\text{rand}}_{e} - \mathrm{PDR}^{\text{atk}}_{e} \Big),
  \label{eq:adversarial_gap}
\end{equation}

\noindent reported with a paired $95\%$ confidence interval,
$\Delta^{\text{adv}} \pm t_{0.975,\,n-1}\, s_d / \sqrt{n}$, where $s_d$ is the
standard deviation of the per-episode differences. Differencing within episode pairs
removes the episode-to-episode variance that otherwise dominates the comparison. The
raw drop with respect to the clean arm is reported alongside it, but is not used as
the measure of attack strength.

\subsubsection*{Decisions versus outcomes}

Two families of metric are recorded and reported separately. \emph{Decision} metrics
count how much routing behaviour the attack changed, principally the action flip
rate, the fraction of per-agent, per-destination $\arg\max$ path choices that differ
between the clean and adversarial observations at the same time step,
\begin{equation}
  \mathrm{flip}
  = \frac{1}{|\mathcal{C}|\, n_d\, T_{\text{ep}}}
    \sum_{t}\sum_{i \in \mathcal{C}}\sum_{d}
    \mathbb{1}\Big[
      \arg\max_k \pi_{\theta_i}\big(\hat{s}_{i,t}\big)_{d,k}
      \neq
      \arg\max_k \pi_{\theta_i}\big(s_{i,t}^{\text{local}}\big)_{d,k}
    \Big].
  \label{eq:flip_rate}
\end{equation}

\emph{Outcome} metrics measure delivered traffic, principally end-to-end PDR. The
distinction is essential in a topology with $K$-path redundancy, where a high flip
rate accompanied by an unchanged PDR indicates that the network absorbed the
redirected flows rather than that the attack succeeded. Reporting a flip rate as
though it were a measure of damage would conflate the two.
